\documentclass[11pt,a4paper]{article}

\usepackage[margin=1in]{geometry}
\usepackage{times}
\usepackage{booktabs}
\usepackage{xcolor}
\usepackage{listings}
\usepackage{enumitem}
\usepackage[hidelinks]{hyperref}
\usepackage{titlesec}

\definecolor{codebg}{RGB}{245,247,249}
\definecolor{codekw}{RGB}{27,94,140}
\definecolor{codecm}{RGB}{110,120,130}

\lstset{
  basicstyle=\ttfamily\small,
  backgroundcolor=\color{codebg},
  keywordstyle=\color{codekw}\bfseries,
  commentstyle=\color{codecm}\itshape,
  language=Python,
  breaklines=true,
  frame=single,
  rulecolor=\color{codebg},
  showstringspaces=false,
  columns=fullflexible,
  xleftmargin=6pt,
  aboveskip=8pt,belowskip=8pt
}

\newcommand{\code}[1]{\texttt{\detokenize{#1}}}

\titleformat{\section}{\large\bfseries}{\thesection}{0.6em}{}
\titleformat{\subsection}{\normalsize\bfseries}{\thesubsection}{0.6em}{}

\title{\textbf{DhakaSim: Java-to-Python Reimplementation and\\
Interface/Architecture Improvements}\\[4pt]
\large Progress Report}
\author{Mehzad Galib}
\date{\today}

\begin{document}
\maketitle

\begin{abstract}
\noindent
DhakaSim is an event-based microscopic traffic simulator for the non-lane-based,
heterogeneous traffic characteristic of cities such as Dhaka. This report
documents the work completed so far: the reimplementation of the original Java
codebase as a dependency-free Python package validated for numeric fidelity
against the Java reference, and a set of improvements to the graphical interface,
the output architecture, and the run reporting. Together these turn a faithful
but bare research port into a self-documenting, easier-to-use tool.
\end{abstract}

\section{Introduction}
Traffic in Dhaka is non-lane-based and heterogeneous: motorised and non-motorised
vehicles share the same right-of-way without lane discipline. DhakaSim models this
by dividing each road into thin lateral \emph{strips} rather than lanes; a vehicle
occupies several strips and may move to any free strip. This report covers two
strands of completed work: (i)~the Java-to-Python reimplementation that forms the
foundation (Section~\ref{sec:port}); and (ii)~the interface, architecture and
reporting improvements built on top of it (Sections~\ref{sec:gui}
and~\ref{sec:arch}). Section~\ref{sec:summary} lists the files changed.

\section{The Python Reimplementation of DhakaSim}
\label{sec:port}
The simulator was reimplemented from its original Java form as a self-contained
Python package (Python~3.10+, standard library only; the GUI uses \code{tkinter}).
The goal was behavioural equivalence with the Java reference rather than idiomatic
rewriting.

\subsection{Translation approach}
The translation follows a one-module-per-class correspondence: each original Java
class becomes a Python module of the same responsibility---for example
\code{vehicle.py}, \code{node.py}, \code{link.py}, \code{segment.py},
\code{strip.py}, \code{processor.py}, \code{signal.py}, \code{pedestrian.py} and
\code{intersection_strip.py}. The package reads its network, demand and parameter
files from \code{input/} and writes results to \code{statistics/}, mirroring the
original layout.

\subsection{Numeric-fidelity layer}
Python's default arithmetic differs from Java's in ways that would shift
simulation results, so a dedicated compatibility layer, \code{javacompat.py},
reproduces the Java semantics that matter:
\begin{itemize}[nosep,leftmargin=1.4em]
  \item \code{min}/\code{max} that propagate \code{NaN} (Gipps braking produces
        \code{NaN} routinely, and mishandling it shifts every speed statistic);
  \item half-up rounding rather than banker's rounding;
  \item integer casts that truncate toward zero;
  \item division by zero and \code{sqrt} of a negative yielding
        \code{Infinity}/\code{NaN} instead of raising;
  \item the sign convention of \texttt{\%} on negative operands;
  \item \texttt{"\%.3f"}-style formatting that rounds the shortest round-tripping
        decimal half-up;
  \item a seed-faithful reimplementation of \code{java.util.Random}.
\end{itemize}
These behaviours are pinned by a regression test, \code{tests/test_javacompat.py}.

\subsection{Validation against the Java reference}
With a matched random seed, the Python build reproduces byte-identical
\code{statistics/csv/*.csv} files and identical printed metrics relative to the
Java implementation, across all thirteen car-following models, all four
lane-changing models, the pedestrian and roadside-object paths (including the
per-step accident log), the route/demand generator output, and the drawing
geometry written to \code{trace.txt}. Unseeded, repeated runs agree within
statistical noise.

\subsection{Testing and execution}
The build was exercised end-to-end in headless mode: it completes a run, prints
the summary metrics, and produces the full set of per-type and per-route output
files. The numeric regression test passes, confirming the Java-compatible
arithmetic has not drifted.

\section{GUI and Usability Improvements}
\label{sec:gui}
A set of changes make the interface understandable to first-time users without
reading the source.

\subsection{Self-documenting option panel}
Every field on the ``Start Simulation'' screen now carries its unit in the label
and a plain-language explanation beside it---for example, the simulation end time
in seconds, the simulation speed in milliseconds per frame, strip width in metres
with the non-lane explanation, and maximum speed in km/h---along with a note of
what the random seed and trace mode do. A subtitle explains that a report opens
automatically when the run ends.

\subsection{Named nodes}
Network nodes can now be labelled by name rather than by numeric id. Names are read
from an optional \code{input/node_names.txt} file (one \code{id name} pair per
line), so any network can supply its own; nodes without a name fall back to their
id. A latent bug was fixed in passing: junction labels, which had been drawn at the
coordinate origin because junction centres are stored as $(0,0)$, are now placed at
the true junction position, recovered as the mean of the link endpoints that meet
there.

\subsection{Units in terminal output}
The summary metrics printed at the end of a run now carry units: speeds are
reported in km/h and waiting times in seconds, matching the report.

\subsection{Per-type vehicle colours}
Vehicles were previously drawn in random colours, making them indistinguishable by
type. A fixed colour per type is now defined once and used by both the animation
and the report, with cars, buses and trucks sharing red, purple and grey families
respectively and CNGs drawn green. The original random draw is preserved so runs
remain seed-identical.
\begin{lstlisting}
VEHICLE_TYPE_COLORS = (
    (0,160,160),  (230,130,20), (140,90,40),  (40,110,220),   # bicycle..motorbike
    (210,50,50),  (170,30,60),  (240,105,105),                # car x3
    (30,160,70),  (140,60,190), (200,70,170),                 # CNG, bus x2
    (100,100,100),(55,55,55),   (0,0,0),                      # truck x2, pedestrian
)
\end{lstlisting}

\subsection{Collapsible in-GUI legend}
A small legend panel on the right of the simulation window maps each vehicle colour
to its type, drawn from the same colour table as the vehicles. It is collapsible via
a $-$/$+$ toggle so it can be tucked away on narrow screens.

\section{Architecture and Reporting Improvements}
\label{sec:arch}

\subsection{Automated HTML report}
A new module, \code{report.py}, writes a self-contained, dependency-free HTML
report at the end of every run. It presents headline metric cards (overall,
motorised and non-motorised speed and waiting), a ``what was simulated''
configuration table, a per-vehicle-type results table (generated, trips completed,
speed, waiting~\%, trip time, fuel), two SVG bar charts (vehicles generated by type
and average speed by type), a colour legend, and a glossary defining every term. In
the GUI, a dialog offers to open the report when the run finishes. The summary
values are derived from the in-memory statistics rather than by re-reading files,
so the report is independent of the CSV outputs.

\subsection{Reorganised statistics output}
The output layout was cleaned up so that the \code{statistics/} folder holds only
the human-readable reports. Each run's report is written with a time-stamped
filename, \code{report_YYYYMMDD_HHMMSS.html}, so successive runs are retained like
log files rather than overwriting one another. The raw numeric CSVs (per-type,
per-route, flow and the per-event accident log) are routed into a
\code{statistics/csv/} subfolder, keeping the folder root uncluttered while
preserving the full numeric record. This routing is centralised, so every CSV
writer is redirected in one place.

\subsection{Offline visualisation}
Because the interactive \code{tkinter} window must run on the user's own machine, an
offline renderer was produced that reuses the simulator's own vehicle-drawing
geometry to rasterise frames into a still image and an animated preview. This makes
it possible to inspect a run's behaviour without a live display, and it shares the
per-type colour scheme so the output matches the in-GUI legend.

\section{Summary of Changes}
\label{sec:summary}
Table~\ref{tab:files} lists the source files added or modified in this work.

\begin{table}[h]
\centering
\small
\begin{tabular}{@{}p{0.30\textwidth}p{0.64\textwidth}@{}}
\toprule
\textbf{File} & \textbf{Change} \\
\midrule
\code{dhakasim/report.py}      & New: HTML report generator (metric cards, tables, SVG charts, colour legend, glossary, time-stamped files) \\
\code{dhakasim/gui.py}         & Annotated option panel; named nodes with corrected label placement; report dialog on finish; collapsible colour legend \\
\code{dhakasim/processor.py}   & Per-type vehicle colours; node-name loading; report hook; km/h units in summary output; CSVs routed to \code{statistics/csv/} \\
\code{dhakasim/constants.py}   & Fixed per-type vehicle-colour palette \\
\code{dhakasim/parameters.py}  & Storage for optional node names \\
\code{dhakasim/vehicle.py}     & Accident log written under \code{statistics/csv/} \\
\code{input/node_names.txt}    & New: optional friendly node names \\
\code{README.md}               & Documentation updated for the new output layout, node names and units \\
\bottomrule
\end{tabular}
\caption{Source files added or modified.}
\label{tab:files}
\end{table}

\section{Conclusion}
\label{sec:conclusion}
The simulator now stands as a validated Python reimplementation of DhakaSim that is
considerably easier to understand and operate: named nodes, unit-labelled output,
type-coloured vehicles and an in-GUI legend on the interface side, and an automatic,
self-explaining per-run HTML report plus a tidy, time-stamped output layout on the
architecture side. The numeric-fidelity guarantees of the original port are
preserved throughout---all colour and reporting changes leave the random stream and
computed statistics unchanged.

\begin{thebibliography}{9}
\small
\bibitem{dev2014} H.\ Dev \emph{et al.}, ``DhakaSim: A Tool for Simulating Traffic
Mobility in Dhaka City,'' \emph{MoHCI}, 2014.
\bibitem{mushfiq2024} M.\ M.\ Mushfiq \emph{et al.}, ````To Lane or Not to
Lane?''---Comparing On-Road Experiences in Developing and Developed Countries Using
a New Simulator RoadBird,'' \emph{IEEE Trans.\ Intelligent Transportation Systems},
vol.~25, no.~8, pp.~8486--8498, 2024.
\end{thebibliography}

\end{document}
